\documentclass[10pt, aspectratio=169]{beamer}

%%\usetheme{metropolis}
\usetheme[progressbar=frametitle]{metropolis}
\usepackage{appendixnumberbeamer}
\usepackage{listings}
\usepackage{proof}
\usepackage{biblatex}
\usepackage{xspace}
\usepackage{tikz-cd}
\newcommand{\themename}{\textbf{\textsc{metropolis}}\xspace}
\newcommand{\U}{\mathcal{U}}
\newcommand{\deq}{\equiv}
\newcommand{\Prod}[2]{\prod_{(#1)}{#2}}
\newcommand{\Sum}[2]{\sum_{(#1)}{#2}}
\newcommand{\W}[2]{W_{(#1)}{#2}}
\newcommand{\0}{\mathbf{0}}
\newcommand{\1}{\mathbf{1}}
\newcommand{\2}{\mathbf{2}}
\newcommand{\N}{\mathbb{N}}
\newcommand{\Z}{\mathbb{Z}}
\newcommand{\Id}{\mathsf{Id}}
\newcommand{\ind}{\mathbf{ind}}
\newcommand{\refl}{\mathsf{refl}}
\newcommand{\Sph}{\mathbb{S}}
\newcommand{\inl}{\mathbf{ind}}

\newtheorem{axiom}{Axiom}

\AtBeginSection[]{
  \begin{frame}
    \sectionpage
  \end{frame}
}

\addbibresource{references.bib}

\title{Type Theory in a Lean Context}
\subtitle{An Alternative Foundation to Mathematics}
\date{October 12, 2026}
\author{Naveen Maurya, IISc}
\institute{Feynman Lecture Hall, ICTS}

\begin{document}

\maketitle

\begin{frame}{Table of contents}
  \setbeamertemplate{section in toc}[sections numbered]
  \tableofcontents[hideallsubsections]
\end{frame}

\section{Foudational Systems} 
 
\subsection{What's wrong with Set Theory?}

\begin{frame}{What is a Foundational System?}
  Set Theory is not a ``complete" foundation mathematics.
  \begin{itemize}
    \item Set theory has a problem about interpreting \emph{types} of mathematical objects.\\
    For example, $\{\{\}, \{\{\}\}\}$ could mean either $2$ or the ordered pair $(0,0)$.
    \item Set theory doesn't center the type of organization mathematicians care about. \\
    For example, definitions of a \emph{group} are closely related
    but that is not reflected in Set Theory.
  \end{itemize}
  %Other problems:
  %\begin{itemize}
  %  \item Set theory is not very suitable for \emph{constructive mathematics}.
  %  \item Hard to verify!!! (by a computer)
  %\end{itemize}
\end{frame}


\subsection{What is a Foundational System?}

\begin{frame}{What is a Foundational System?}
    A foundational system is an assumed framework to build on top of. \\
    \[
    \begin{tikzcd}
    ZFC 
    \end{tikzcd}
    \]
\end{frame}

\begin{frame}{What is a Foundational System?}
  Concretely, a Foundation System consists of
  \begin{itemize}
    \item A \textbf{Logical/Deductive system} specifying \emph{rules of inference} that dictate which inferences are
      allowed from existing context.
    \item An \textbf{Axiomatic system}, a schematic collection of \emph{axioms}, statements in the language of the deductive system,
      taken as true by assumption.
  \end{itemize}
  %Usually in set-theoretic foundations, the former is \emph{first-order logic} and the latter is \emph{ZF or ZFC Set Theory}.
  %Type Theory combines these two layers into the same framework.
\end{frame}

\section{(Martin-L\"of) Type Theory}

\subsection{Types as Mathematical Building Blocks}

\begin{frame}{Types, Terms and Universes}
  \emph{Types} and \emph{Terms} are primitive objects.
  \begin{itemize}
    \item Every term $a$ is of some (unique) type, say $\alpha$. We denote this $a : \alpha$. 
  \end{itemize}
\end{frame}

\begin{frame}{Universes}
  Type Theory postulates a chain of cumulative universes in the style of Russell:
  \[
    \U_0, \U_1, \U_2, \U_3, \ldots
  \]
  satisfying:
  \begin{itemize}
    \item $\U_i : \U_{i+1}$ for each $i$
    \item If $A : \U_i$, then $A : \U_{i+1}$
  \end{itemize}
  Each type is a term of some universe.
  Each universe is also a term of universes above it.
\end{frame}

\section{Logic in Type Theory}

\begin{frame}{Embedding a Logic through the Construction Material}
Type Theory repurposes the available bricks for building math to also encode logic!
\begin{itemize}
  \item Certain constructions and manipulations (type formers) of \emph{bricks} are available in the system.
\end{itemize}
Certain types, called \emph{mere propositions} may be used to capture logical propositions.\\
One may intepret terms of these types as \emph{proofs} to the proposition. \\
Hence, proofs of statements are themselves first-class mathematical objects.
\end{frame}

\begin{frame}{Type Formers}
  %\setlength{\tabcolsep}{15pt}     % increase column spacing
  \centering
  \begin{tabular}{c c c c}
  \hline
  \textbf{Type Former} & \textbf{Assumes} & \textbf{Notation} & \textbf{Logical Analogue} \\
  \hline
  Function Type             & $A, B : \U$           & $A \to B$             & $A \implies B$ \\
  Dependent Function Type   & $A : \U, B : A \to \U $ & $\Prod{x:A}{B(x)}$  & $\forall x \in A, B(x)$ \\
  Dependent Pair Type       & $A : \U, B : A \to \U$  & $\Sum{x:A}{B(x)}$   & $\exists x \in A, B(x)$ \\
  Product Type              & $A, B : \U$           & $A \times B$             & $A \land B$ \\
  Coproduct Type            & $A, B : \U$           & $A + B$             & $A \lor B$ \\
  Empty Type                & & $\0$                & $\mathtt{False}$ \\
  Unit Type                 & & $\1$                & $\mathtt{True}$ \\
  \hline
  \end{tabular}
\end{frame}

\begin{frame}{Constructivity of Type Theory}
  Type Theory is a constructivist foundational system!
  \begin{itemize}
    \item A function $f : A \to B$ needs to have an alogorithm to effectively output an element $f(a) : B$ when given input $a : A$.
    \item $p \lor q$ requires either a proof of $p$ or a proof of $q$.
    \item A proposition $\exists a : A, B(a)$ needs a an explicit $a : A$ such that $B(a)$ as a certificate.
  \end{itemize}
  Curry-Howard Correspondence b/w proofs and programs! \\
  Type Theory requires proofs to be \emph{computationally meaningful}.
\end{frame}

\section{Mathematics in Type Theory}

\subsection{Inductive Types}

\begin{frame}{Inductive Types}
  Natural numbers are the \emph{model mathematical object} in Type Theory!
  
  Remember Peano Axioms for $\N$,
  \begin{itemize}
    \item There is a distinguished element $0 \in \N$. 
    \item There is an injective function $\mathtt{succ} : \N \to \N$, 
    \item There is no $n \in \N$ such that $\mathtt{succ}(n) = 0$.
    \item The \emph{induction principle}, i.e. to show a predicate $P$ on $\N$ holds for all $n \in \N$,
      it suffices to prove $P(0)$ holds, and that for all $n \in \N, P(n) \implies P(\mathtt{succ} (n))$.
  \end{itemize}
\end{frame}

\begin{frame}{Inductive Types}
  There's a somewhat obvious-looking property one can prove:
  \begin{theorem}[Classification of Natural Numbers]
    Any $n \in \N$ is either $0$ or $\mathtt{succ}^m(0)$ for some $m \in \N$.
  \end{theorem}
  \begin{proof}
    Let $P(n)$ be the statement ``$n = 0$ or there exists some $m \in \N$ such that $\mathtt{succ}^m(0) = n.$" \\
    By the induction principle, it suffices to show $P(0)$ and $\forall n \in \N, P(n) \implies P(\mathtt{succ}(n))$.

    $P(0)$ is clearly true, and for some $n \in \N$, suppose $n = \mathtt{succ}^m(0)$, then note $\mathtt{succ}(n) = \mathtt{succ}(\mathtt{succ}(m)(0)) = \mathtt{succ}^{\mathtt{succ}(m)}(0)$.
  \end{proof}
    %Let $P(n)$ be the statement ``$n = 0$ or there exists some $m \in \N$ such that $\mathtt{succ}(m) = n.$" \\
    %By the induction principle, it suffices to show $P(0)$ and $\forall n \in \N, P(n) \implies P(\mathtt{succ}(n))$.

    %$P(0)$ is clearly true, and for any $n \in \N$, there is $m := n \in \N$ such that $\mathtt{succ}(n) = \mathtt{succ}(m)$.
    %This completes the proof.

\end{frame}

\begin{frame}{Inductive Types}
  Here's a \emph{universal property} about $\N$ wrt maps out of $\N$. Given $E : \N \to \U$,
    \[\Prod{n : \N}{E(n)} \simeq \left(E(0) \times \Prod{n : \N}{E(n) \to E(\mathtt{succ}(n))} \right)\]
\end{frame}

\begin{frame}{Inductive Types}
  There is a general schema for constructing new types, requiring:
  \begin{itemize}
    \item \textbf{Formation rules} to introduce an inductive type into context. 
    \item \textbf{Constructors/introduction rules} to introudce terms of the type.

    \item \textbf{Elimination rules} to consume terms of a type. \\
      (Usually by an \emph{induction} principle which defines dependent functions on the type)

    \item \textbf{Computation rules} to specify action of elimination on the constructors.

  \end{itemize}
  These rules specify the behaviour of an inductive type by dictating the inference rules for the type 
  in the foundational system.
\end{frame}

\subsection{Examples}

\begin{frame}{Examples}
  \textbf{Natural Numbers Type ($\N$)}
  The natural numbers are an inductive type, given by the constructors 
  \begin{itemize}
    \item $\mathtt{zero} : \N$ 
    \item $\mathtt{succ} : \N \to \N$
  \end{itemize}
  and the induction principle that for $E : \N \to \U$, to prove $\Prod{n : \N}{E(n)}$ is equivalent to proving
  $E(0)$ and $\Prod{n : \N} E(n) \to E(\mathtt{succ}(n))$.
  Formally, 
  \[
  \ind_\N : \Prod{E : \N \to \U} (E(0) \times \Prod{n : \N} (E(n) \to E(\mathtt{succ}(n)))) \to \Prod{n:\N}{E(n)}
  \]
\end{frame}

\begin{frame}{Examples}
  \textbf{The Empty Type ($\0$)}
  The empty type $\0$ is given by no constructors.
  One has the universal function for any $E : \0 \to \U$, one has the universal function $\Prod{z : \0}{E(z)}$.
  Formally, 
  \[
  \ind_0 : \Prod{E : \0 \to \U}{\Prod{z : \0}{E(z)}}
  \]
\end{frame} 
\begin{frame}{Examples}
 \textbf{The Unit Type ($\1$)}
  The unit type $\1$ is given by one constructor,
  \begin{itemize}
    \item $* : \1$.
  \end{itemize}
  and the induction principle that for $E : \1 \to \U$, to prove $\Prod{z : \1}{E(z)}$, it suffices to prove $E(*)$.
  Formally, 
  \[
  \ind_1 : \Prod{E : \1 \to \U}{E(*) \to \Prod{z : \1}{E(z)}}
  \]
\end{frame}

\begin{frame}{Examples}
  \textbf{The Coproduct Type ($A + B$)}
  The coproduct type of $A$ and $B$ is an inductive type, given by the constructors 
  \begin{itemize}
    \item $\mathtt{inl} : A \to A + B$ 
    \item $\mathtt{inr} : B \to A + B$
  \end{itemize}
  and the induction principle that for $E : A + B \to \U$, to prove $\Prod{z : A + B}{E(z)}$ is equivalent to proving 
  $E(\mathtt{inl} (a))$ for all $a : A$ and $E(\mathtt{inr} (b))$ for all $b : B$.
  Formally, 
  \[
  \ind_{A+B} : \Prod{E : A+B \to \U}{\left( \Prod{a : A}{E(\mathtt{inl}(a))} \times \Prod{a : A}{E(\mathtt{inl}(a))} \to \Prod{z : A + B}{E(z)} \right) }
  \]
\end{frame}

%\begin{frame}{The General Construction}
%  \textbf{W-Types}
%  In general, given $A : \U$, the type of \emph{constructors}, and type families $B : A \to \U$, we are allowed an inductive type
%  \begin{itemize}
%  \item $\mathtt{sup} : \left( \Prod{a : A}{B(a)} \to \W{a:A}{B} \right) \to \W{a:A}{B(a)} $
%  \end{itemize}
%  and the induction principle that for $E : \W{a:A}{B(a)} \to \U$, to prove $\Prod{z : A + B}{E(z)}$ is equivalent to proving 
%  $E(\mathtt{inl} (a))$ for all $a : A$ and $E(\mathtt{inr} (b))$ for all $b : B$.
%  Formally, 
%  \[
%  \ind_{\W{a:A}{B(a)}} : \Prod{E : A+B \to \U}{\left( \Prod{a : A}{E(\mathtt{inl}(a))} \times \Prod{a : A}{E(\mathtt{inl}(a))} \to \Prod{z : A + B}{E(z)} \right) }
%  \]
%\end{frame}
%

%\begin{frame}{The General Construction}
%  \textbf{W-Types}
%  In general, given $A : \U$, the type of \emph{constructors}, and type families $B : A \to \U$, we are allowed an inductive type
%  \begin{itemize}
%  \item $sup : \left( \Prod{a : A}{B(a)} \to \W{a:A}{B} \right) \to \W{a:A}{B(a)} $
%  \end{itemize}
%  and the induction principle that for $E : \W{a:A}{B(a)} \to \U$, to prove $\Prod{z : A + B}{E(z)}$ is equivalent to proving 
%  $E(\mathtt{inl} (a))$ for all $a : A$ and $E(\mathtt{inr} (b))$ for all $b : B$.
%  Formally, 
%  \[
%  \ind_{\W{a:A}{B(a)}} : \Prod{E : A+B \to \U}{\left( \Prod{a : A}{E(\mathtt{inl}(a))} \times \Prod{a : A}{E(\mathtt{inl}(a))} \to \Prod{z : A + B}{E(z)} \right) }
%  \]
%\end{frame}

\begin{frame}{A Custom Inductive Type}
  \textbf{Binary Trees}
  Given a type $\mathtt{A} : \U$, the type of binary trees over $\mathtt{A}$ is $\mathtt{BinTree(A)}$ given by the constructors
  \begin{itemize}
  \item $\mathtt{leaf} : \mathtt{A} \to \mathtt{BinTree(A)}$
  \item $\mathtt{node} : \mathtt{BinTree(A)} \to \mathtt{BinTree(A)} \to \mathtt{BinTree(A)}$
  \end{itemize}
  with the induction principle, that for $E : A \to \U$, to prove $\Prod{z : \mathtt{BinTree(A)}}{E(z)}$,
  it suffices to prove $\Prod{a : A}{E(\mathtt{leaf(a)})}$ and $\Prod{t_1, t_2 : \mathtt{BinTree(A)}}{E(\mathtt{node}(t_1)(t_2))}.$
\end{frame}

\begin{frame}{Structural Constructions}
  \textbf{Magma}
  A \emph{magma} is a type $M$ with a binary operation $\cdot : M \times M \to M$. The type $Magma$ is given by
  \[
  Magma :\deq \Sum{M : \U}{M \times M \to M}
  \]
  \textbf{Semigroup}
  A \emph{semigroup} is a magma with the operation being associative
  \[
  Semigroup :\deq \Sum{M : \U}{\Sum{\cdot : M \times M \to M}{\Prod{a, b, c : M}{(a \cdot b) \cdot c = a \cdot (b \cdot c)}}}
  \]
\end{frame}



\subsection{Identity Types}

\begin{frame}{Identity Types}
  What does it mean for two mathematical objects to be equal? \\
  In (intensive) type theory, there are two types of equalities:
  \begin{itemize}
    \item \textbf{Definitional Equality} Equality within the logical framework.
    \item \textbf{Propositional Equality} A genuine equality of defined mathematical objects.
  \end{itemize}
  \pause
  \textbf{Remark.}
  An equality $p : A =_\U B$ in a universe gives a transport $p_* : E(A) \to E(B)$ for any $E : A \to \U$.
\end{frame}

\section{Conclusion}

\begin{frame}{Conclusion}
\end{frame}

\begin{frame}{References}
\printbibliography
\nocite{*}
\end{frame}

\begin{frame}
    \begin{center}
        \Huge{Thank You!}
    \end{center}
\end{frame}
\end{document}
